% =============================================================================
%  dd_solver_simple.tex — the mathematics of cpp/dd_solver_simple.hpp
%
%  A self-contained description of the readable domain-decomposition
%  (arrowhead) linear solver for IPOPT: the arrowhead permutation, the Schur
%  complement solve, the Haynsworth inertia, the forward-only + sparse-RHS +
%  row-contraction formation of the local Schur blocks, the FETI-DP-style
%  peel, ASd-preconditioned conjugate gradients on the interface, and the
%  predicted inertia In(S) = In(T).
%
%  Build:  pdflatex dd_solver_simple.tex  (twice, for references)
% =============================================================================
\documentclass[11pt]{article}

\usepackage[margin=1.1in]{geometry}
\usepackage{amsmath,amssymb,amsthm}
\usepackage{booktabs}
\usepackage{microtype}
\usepackage{xcolor}
\usepackage[colorlinks=true,linkcolor=blue!50!black,citecolor=blue!50!black,urlcolor=blue!50!black]{hyperref}
\usepackage{algorithm}
\usepackage{algpseudocode}

\newtheorem{theorem}{Theorem}
\newtheorem{proposition}{Proposition}
\newtheorem{remark}{Remark}

% ---- notation ---------------------------------------------------------------
\newcommand{\In}{\operatorname{In}}
\newcommand{\diag}{\operatorname{diag}}
\newcommand{\R}{\mathbb{R}}
\newcommand{\code}[1]{\texttt{#1}}
\newcommand{\Sff}{S_{\mathrm{ff}}}
\newcommand{\SfP}{S_{\mathrm{fP}}}
\newcommand{\SPP}{S_{\mathrm{PP}}}

\title{The Arrowhead Domain-Decomposition Linear Solver\\
       of \code{dd\_solver\_simple.hpp}:\\
       Mathematics and Design}
\author{David Villacis}
\date{September 2026}

\begin{document}
\maketitle

\begin{abstract}
\noindent
\code{dd\_solver\_simple.hpp} is a complete, self-contained linear solver for
IPOPT built only on Eigen: no HSL (MA57/MA97), no MUMPS.  Given a partition of
the KKT unknowns into $K$ subdomains plus a shared border, it permutes every
Newton system into bordered block-diagonal (``arrowhead'') form, factorizes
the subdomain blocks independently by sparse $LDL^{\mathsf T}$, and solves the
coupling (Schur-complement) system by preconditioned conjugate gradients ---
\emph{without ever assembling or factorizing the Schur complement}.  The
inertia that IPOPT requires is delivered by two applications of the Haynsworth
additivity identity: once across the subdomain blocks, and once more across a
small ``peel'' split of the interface, so that the only globally coupled
computation left in the algorithm is the inner product reductions of the
Krylov iteration.  This note records every piece of mathematics the
implementation relies on, the failure semantics that make the method safe
inside IPOPT's regularization loop, and the measurements that selected this
particular configuration from the larger design space explored in the
production solver \code{dd\_solver.hpp}.
\end{abstract}

\tableofcontents

% =============================================================================
\section{The contract with IPOPT}
\label{sec:contract}

At every iteration of its interior-point method, IPOPT
\cite{waechter2006implementation} must solve the primal--dual Newton system.
After the barrier terms and its own regularization are folded in, this is a
symmetric indefinite \emph{augmented KKT system}
\begin{equation}
  A\,\Delta z \;=\; r,
  \qquad
  A \;=\;
  \begin{bmatrix}
    H + \Sigma + \delta_w I & J^{\mathsf T}\\[2pt]
    J & -\delta_c I
  \end{bmatrix}
  \in \R^{n_A\times n_A},
  \label{eq:kkt}
\end{equation}
where $H$ is the Lagrangian Hessian, $\Sigma$ the barrier diagonal
(entries scaling like $z^2/\mu$, which is what drives $\|A\|$ to
$\sim\!10^{18}$ near convergence), $J$ the constraint Jacobian, and
$\delta_w,\delta_c \ge 0$ IPOPT's inertia-correction and dual-regularization
shifts, \emph{already applied} by the time the matrix reaches the linear
solver.  The matrix is handed over in lower-triangle triplet form
$(i,j,v)$, one-based, possibly with duplicate entries that must be summed.

The linear solver must return two things:
\begin{enumerate}
\item a solution of $A\,\Delta z = r$ (possibly for several right-hand
      sides against one factorization), and
\item the \emph{inertia} of $A$ --- specifically $n_-(A)$, the number of
      negative eigenvalues.
\end{enumerate}
Item 2 is not optional.  Writing
$\In(A) = \bigl(n_+(A),\,n_-(A),\,n_0(A)\bigr)$
for the numbers of positive, negative, and zero eigenvalues, IPOPT's
inertia-correction loop accepts a step only when
$n_-(A) = m$ (the number of constraint rows) and $n_0(A)=0$; otherwise it
increases $\delta_w$ (and, on singularity, $\delta_c$) and asks for a new
factorization.  A solver that misreports $n_-$ silently corrupts this loop:
the filter line search then fights a curvature defect that does not exist.
A solver that cannot report inertia at all forces IPOPT into its inertia-free
curvature test \cite{chiang2016inertia} --- a mode this solver deliberately
does \emph{not} use (Section~\ref{sec:inertia-pred} explains what it uses
instead, and Section~\ref{sec:measured} records why).

Three status outcomes cover all failure semantics:
\begin{center}
\begin{tabular}{ll}
\toprule
\code{SYMSOLVER\_SUCCESS} & factorization and solves usable\\
\code{SYMSOLVER\_SINGULAR} & factorization unusable; IPOPT raises
  $\delta_w,\delta_c$ and retries\\
\code{SYMSOLVER\_WRONG\_INERTIA} & $n_-\ne m$; IPOPT raises $\delta_w$
  and retries\\
\bottomrule
\end{tabular}
\end{center}
Both retry paths are \emph{regularization} paths.  A guiding principle of the
implementation is therefore: report \code{SINGULAR} only for conditions that a
larger $\delta_w$ can actually cure (a broken-down factorization, an
undeterminable inertia), and never for conditions it cannot (an iterative
solve that failed to converge --- see Section~\ref{sec:refine}).

IPOPT owns everything else: the filter line search, the $\mu$-homotopy, the
restoration phase.  The solver owns only the linear algebra of
\eqref{eq:kkt}.

% =============================================================================
\section{The arrowhead form}
\label{sec:arrowhead}

The caller provides an \emph{owner map}: for every KKT index
$i \in \{1,\dots,n_A\}$ either a subdomain label $k \in \{1,\dots,K\}$ or the
marker ``border'' for a \emph{complicating} unknown shared by two or more
subdomains.  (For the image-reconstruction problems this solver was built for,
the map comes from a tile or strip partition of the pixel grid; the solver
itself is agnostic --- it only requires that the map be consistent, in the
sense of \eqref{eq:noleak} below.)

Let $\Pi$ be the permutation that orders all interior unknowns first, grouped
by subdomain, and the $p$ border unknowns last.  Then
\begin{equation}
  \Pi A \Pi^{\mathsf T}
  \;=\;
  \begin{bmatrix}
    W_1 &        &        & B_1^{\mathsf T}\\
        & \ddots &        & \vdots\\
        &        & W_K    & B_K^{\mathsf T}\\
    B_1 & \cdots & B_K    & C
  \end{bmatrix},
  \label{eq:arrow}
\end{equation}
with $W_k \in \R^{n_k\times n_k}$ the interior block of subdomain $k$,
$B_k \in \R^{p\times n_k}$ its coupling to the border, and
$C \in \R^{p\times p}$ the border--border block.  The block-diagonal shape
requires that no triplet couple two \emph{different} subdomains directly:
\begin{equation}
  \text{owner}(i) = k_1 \ \wedge\ \text{owner}(j) = k_2 \ \wedge\ k_1\ne k_2
  \;\Longrightarrow\; A_{ij} = 0 ,
  \label{eq:noleak}
\end{equation}
which the implementation verifies while routing the triplets and rejects as a
``partition leak'' otherwise.

Two remarks on \eqref{eq:arrow}.  First, it is a \emph{pure permutation}: no
unknowns are duplicated and no linking constraints are introduced (in contrast
to duplicate-and-link reformulations), so $C$ is genuinely nonzero and the
$B_k$ carry real Jacobian and Hessian entries.  Second, each $B_k$ is
extremely sparse: only the border unknowns geometrically adjacent to subdomain
$k$ appear in it.  Writing $N_k \subseteq \{1,\dots,p\}$ for the set of
border positions subdomain $k$ actually touches, with $p_k = |N_k|$, the
matrix $B_k$ has nonzero rows only on $N_k$.  We use the same symbol for the
selection operator $N_k \in \R^{p_k \times p}$ picking out those positions.

\subsection{Block elimination}

Assume for the moment every $W_k$ is nonsingular (Section~\ref{sec:ldlt}
handles failure).  Eliminating the interior blocks from \eqref{eq:arrow}
produces the \emph{interface} (Schur complement) matrix
\begin{equation}
  S \;=\; C \;-\; \sum_{k=1}^{K} B_k W_k^{-1} B_k^{\mathsf T}
  \;\;\in\R^{p\times p},
  \label{eq:schur}
\end{equation}
and the three-stage solve for a right-hand side split as
$r = (r_1,\dots,r_K,\,r_y)$:
\begin{align}
  r_S &= r_y - \sum_{k=1}^{K} B_k W_k^{-1} r_k ,
  \label{eq:solve-a}\\
  \Delta y &= S^{-1} r_S ,
  \label{eq:solve-b}\\
  \Delta x_k &= W_k^{-1}\bigl(r_k - B_k^{\mathsf T}\,\Delta y\bigr),
  \qquad k = 1,\dots,K.
  \label{eq:solve-c}
\end{align}
Everything indexed by $k$ is independent across $k$ --- that is the point of
the decomposition.  Only \eqref{eq:solve-b} is globally coupled, and the
design question of this whole solver is how to perform
\eqref{eq:solve-b} \emph{without a global factorization}.

\subsection{Inertia by Haynsworth additivity}
\label{sec:haynsworth}

\begin{theorem}[Haynsworth inertia additivity \cite{haynsworth1968}]
\label{thm:haynsworth}
Let
$M = \begin{bmatrix} E & F^{\mathsf T}\\ F & G\end{bmatrix}$
be symmetric with $E$ nonsingular.  Then
\[
  \In(M) \;=\; \In(E) \;+\; \In\!\bigl(G - F E^{-1} F^{\mathsf T}\bigr).
\]
\end{theorem}

Applied to \eqref{eq:arrow} with
$E = \diag(W_1,\dots,W_K)$ and Schur complement \eqref{eq:schur},
\begin{equation}
  \In(A) \;=\; \sum_{k=1}^{K} \In(W_k) \;+\; \In(S).
  \label{eq:haynsworth}
\end{equation}
The first term is \emph{free}: each $\In(W_k)$ is read off the pivot signs of
a factorization that must happen anyway, locally and in parallel
(Section~\ref{sec:ldlt}).  The second term, $\In(S)$, is the one quantity
that would classically force assembling and factorizing $S$ --- the single
serial step in an otherwise embarrassingly parallel scheme.
Section~\ref{sec:inertia-pred} removes it.

\begin{remark}[Why pivot signs, not eigenvalues]
On these problems $\|A\|$ reaches $\sim\!10^{18}$ through the barrier terms
$\Sigma$, so the \emph{signs of small eigenvalues} are rounding noise and any
eigenvalue-counting approach is hopeless.  The pivot signs of an
$LDL^{\mathsf T}$ factorization, by Sylvester's law of inertia, count
$n_\pm$ exactly regardless of scaling: $\In(A) = \In(D)$ whenever
$A = LDL^{\mathsf T}$ with $L$ unit triangular.  Every inertia statement in
this solver ultimately reduces to a pivot-sign or exact-eigendecomposition
count of a \emph{well-scaled small} matrix.
\end{remark}

% =============================================================================
\section{Subdomain factorization: unpivoted sparse \texorpdfstring{$LDL^{\mathsf T}$}{LDLT}}
\label{sec:ldlt}

Each $W_k$ is factorized by Eigen's \code{SimplicialLDLT}:
\begin{equation}
  W_k \;=\; P_k^{-1}\, L_k D_k L_k^{\mathsf T}\, P_k ,
  \label{eq:ldlt}
\end{equation}
where $P_k$ is a fill-reducing permutation (computed once per sparsity
pattern; the KKT pattern is constant over a whole run, so the symbolic
analysis is reused by every Newton step), $L_k$ is \emph{unit} lower
triangular, and $D_k$ is diagonal.  Eigen applies no scaling, a fact
Section~\ref{sec:sk} depends on.  The inertia is
$\In(W_k) = \In(D_k)$, counted from the signs of the pivots.

\code{SimplicialLDLT} does \emph{not} pivot for stability, and this is the
one real weakness of the Eigen route.  The saving grace is a structural
property of \eqref{eq:kkt}:

\begin{proposition}[Quasidefiniteness \cite{vanderbei1995symmetric}]
If $\delta_w > 0$ makes the $(1,1)$ block of \eqref{eq:kkt} positive definite
and $\delta_c > 0$, then $A$ is symmetric quasidefinite, and \emph{every}
symmetric permutation of $A$ admits an $LDL^{\mathsf T}$ factorization with
nonzero pivots.  The same holds for every $W_k$, which inherits the
saddle-point structure.
\end{proposition}

IPOPT, however, usually offers a matrix with $\delta_w = \delta_c = 0$ first,
and there a zero pivot is possible.  The implementation detects breakdown
exactly (a zero or non-finite entry of $D_k$) and reports
\code{SINGULAR} --- refusing both the solve and the inertia, since neither is
meaningful past a broken pivot.  IPOPT responds with $\delta_w,\delta_c > 0$,
which is \emph{precisely} the regularization that makes the unpivoted
factorization exist.  The failure mode is therefore self-correcting; its cost
is a few extra factorizations per run relative to a pivoting Bunch--Kaufman
code such as MA57, and that is the honest price of dropping HSL.  No
artificial shift is ever added inside the solver: masking a rank deficiency
locally would corrupt the inertia signal \eqref{eq:haynsworth} that drives
IPOPT's $\delta_w$ loop.

% =============================================================================
\section{Forming the local Schur complements}
\label{sec:sk}

The dense $p_k \times p_k$ matrices
\begin{equation}
  S_k \;=\; -\,B_k W_k^{-1} B_k^{\mathsf T}
  \label{eq:sk}
\end{equation}
(so that $S = C + \sum_k N_k^{\mathsf T} S_k N_k$ in selection-operator form)
are the workhorses of everything downstream: the interface operator, the
preconditioner, and the peel all consume them.  Note that $S_k$ is
\emph{mathematically} dense --- entry $(a,b)$ equals
$-\,b_a^{\mathsf T} W_k^{-1} b_b$ with $W_k^{-1}$ full, so generically every
pair of border unknowns touching subdomain $k$ interacts --- hence ``avoid
storing it densely'' is not an available option; the design question is only
how cheaply it can be \emph{formed}.  Formed naively, \eqref{eq:sk} costs
$p_k$ full sparse back-solves (forward, diagonal, backward) with dense
right-hand sides and dominates the entire factorization.  Three exact
observations remove most of that cost.

\subsection{Forward-only formation}

Substituting \eqref{eq:ldlt} into \eqref{eq:sk} and using
$P_k^{-1} = P_k^{\mathsf T}$,
\begin{equation}
  B_k W_k^{-1} B_k^{\mathsf T}
  \;=\; B_k P_k^{-1} L_k^{-\mathsf T} D_k^{-1} L_k^{-1} P_k B_k^{\mathsf T}
  \;=\; Y_k^{\mathsf T} D_k^{-1} Y_k,
  \qquad
  Y_k \;=\; L_k^{-1} P_k B_k^{\mathsf T} \in \R^{n_k\times p_k}.
  \label{eq:forward-only}
\end{equation}
The backward substitution $L_k^{-\mathsf T}$ cancels symbolically against the
$B_k$ that would multiply it back from the left, so \emph{the backward half
of the solve is never performed}, and $B_k$ itself disappears from the
product: $Y_k$ already carries it.  (This identity requires that the
factorization apply no diagonal scaling between the two triangular halves ---
true for Eigen, not for MA57 under MC64 scaling.)

\subsection{Sparse-right-hand-side forward substitution}
\label{sec:sparse-fwd}

Each column of $P_k B_k^{\mathsf T}$ carries only a handful of nonzeros (a
border unknown couples to a few interior unknowns), and a forward
substitution propagates column $j$ of $L_k$ only through the multiplier
$x_j$:\ solving $L_k x = b$ column-by-column in ascending $j$,
\[
  x_i \;\leftarrow\; x_i - (L_k)_{ij}\, x_j
  \quad\text{for all } i > j \text{ with } (L_k)_{ij}\ne 0,
\]
a step that contributes \emph{exactly nothing} when $x_j = 0$.  The columns
of $Y_k$ therefore acquire nonzeros only on the \emph{reach} of the
right-hand-side pattern --- by Gilbert's theorem \cite{gilbert1988}, the
union of elimination-tree paths from the pattern's positions to the root.
Because the KKT pattern is fixed for the whole run, each column's reach is
\emph{static}: it is computed once per sparsity pattern (walking the etree,
which the symbolic analysis already built) and the substitution iterates
exactly that list, ascending --- a valid order since an etree parent always
carries the larger index.  Measured on the target problems the reach covers
$2$--$4\%$ of the columns of $L_k$; a pattern-blind implementation that
instead sweeps all $n_k$ columns testing $x_j = 0$ is algebraically
identical, but at $N=64$ the empty sweep was profiled at more than the
arithmetic it guarded, and the same $O(n_k p_k)$ dense-bookkeeping trap
(scanning and zeroing the mostly-empty $Y_k$ buffer) applies to the
contraction below --- so $Y_k$ is held \emph{compressed on its reach}
throughout.  (Eigen's own sparse-RHS solve path does none of this; it
expands the right-hand side into dense panels.)

\subsection{Row contraction of \texorpdfstring{$Y_k^{\mathsf T} D_k^{-1} Y_k$}{YtDinvY}}
\label{sec:row-contraction}

The remaining cost of \eqref{eq:forward-only} is the contraction
$Y_k^{\mathsf T} D_k^{-1} Y_k$.  As a dense GEMM this costs
$2\,n_k p_k^2$ flops --- almost all of them multiplying the zeros that
Section~\ref{sec:sparse-fwd} just avoided computing (measured: $Y_k$ is
$\approx 88\%$ zeros at $N=32$).  Contracting by \emph{rows} instead,
\begin{equation}
  S_k \;=\; -\sum_{r=1}^{n_k} \frac{1}{(D_k)_{rr}}\; y_r^{\mathsf T} y_r,
  \qquad y_r = \text{the $r$-th row of } Y_k \text{ (sparse)},
  \label{eq:row-contraction}
\end{equation}
costs $\sum_r \operatorname{nnz}(y_r)^2 \le p_k \operatorname{nnz}(Y_k)$
flops, skips only terms that are exactly zero, and produces an exactly
symmetric $S_k$ by construction (the upper triangle is accumulated and
mirrored --- exact symmetry is relied on by the inertia count of
Section~\ref{sec:inertia-pred} and by the preconditioner).  The row
grouping is read off the same static reach lists by a counting sort, so no
dense scan of $Y_k$ ever occurs.  Measured whole-run effect of the two
formation steps (cameraman, exact Hessian, iteration-for-iteration
identical in every cell):
\begin{center}
\begin{tabular}{lcccc}
\toprule
run & dense GEMM & row contraction & + static reach & total\\
\midrule
$N=32$, $3\times 3$ tiles & 4.42\,s & 4.00\,s & 3.79\,s & $1.2\times$\\
$N=64$, $2\times 2$ tiles & 42.8\,s & 20.3\,s & 16.0\,s & $2.7\times$\\
$N=64$, $4$ strips        & 25.7\,s & 10.1\,s & 8.08\,s & $3.2\times$\\
\bottomrule
\end{tabular}
\end{center}

\begin{remark}[Why $S_k$ is nevertheless \emph{kept} dense]
Conjugate gradients only needs the action $p \mapsto S p$, so one could skip
forming $S_k$ entirely: either apply the factored form
$-\,Y_k^{\mathsf T}(D_k^{-1}(Y_k\,p))$ each iteration, or go fully
matrix-free and run one $W_k$ back-solve per iteration with no $Y_k$ at all.
Measured, both lose.  The factored application costs
$\approx 4\operatorname{nnz}(Y_k)$ flops per iteration against $2 p_k^2$ for
the dense matrix--vector product, and $\operatorname{nnz}(Y_k) > p_k^2$ on
most subdomains; at the observed $\sim\!10^3$ CG iterations per
factorization the saved contraction is repaid with interest.  The fully
matrix-free application costs on the order of the union reach of
$\operatorname{nnz}(L_k)$ per iteration --- roughly $10\times$ the dense
product --- and loses by more.  Moreover the ASd preconditioner
(Section~\ref{sec:precond}) consumes the \emph{entries} of $S_k$, not its
action, so the dense block must exist regardless.  The winning combination
is \eqref{eq:row-contraction} on static reaches: sparse arithmetic to
\emph{form} $S_k$, dense arithmetic to \emph{apply} it.
\end{remark}

A runtime referee is kept for all of this: with
\code{DDS\_SCHUR\_CHECK=1} every $S_k$ is recomputed by the naive two-sided
route and compared, per block per factorization (relative error
$\sim\!10^{-12}$ typical, worst $\sim\!10^{-8}$ at the most ill-conditioned
iterates).

% =============================================================================
\section{The interface operator, matrix-free}
\label{sec:matfree}

$S$ itself is \emph{never assembled}.  Its action is evaluated as the same
sum an assembly would have performed,
\begin{equation}
  S\,y \;=\; C\,y \;+\; \sum_{k=1}^{K} N_k^{\mathsf T}\, S_k\, (N_k\, y),
  \label{eq:applyS}
\end{equation}
with every term local to one subdomain.  This is the distributed algorithm,
not a simulation of it: each rank owns its dense $S_k$, and the only
communication per application is one reduction over the border vector.
(The stricter regime in which even the $S_k$ are not kept, and $S y$ is
re-derived through $K$ subdomain back-solves per iteration, is measured in
the production solver as \code{APPLY\_MATFREE}; keeping $S_k$ --- formed once
per Newton step, locally --- is the normal choice.)

The diagonal of $S$, needed by the preconditioner, is likewise assembled
without $S$:
\begin{equation}
  \diag(S) \;=\; \diag(C) \;+\; \sum_{k=1}^{K} N_k^{\mathsf T}\diag(S_k),
  \label{eq:diagS}
\end{equation}
one all-reduce over a $p$-vector in a distributed setting --- which is
exactly why the ASd preconditioner is built around it.

% =============================================================================
\section{The peel: making the interface fit for CG}
\label{sec:peel}

Conjugate gradients requires a symmetric \emph{positive definite} operator,
and $S$ is generally not one.  \emph{Three} kinds of border index are
responsible:

\begin{enumerate}
\item \textbf{Dual border indices.}  On a $k\times k$ tile partition the
  driver promotes the cut-corner dual pairs to the border (otherwise the
  local blocks are structurally rank-deficient).  Empirically the promoted
  duals are \emph{exactly} the negative eigenvalues of $S$: at every recorded
  solve, $n_-(S)$ equalled the promoted count.  (Section
  \ref{sec:why-predictable} explains why.)  So on tile partitions $S$ is
  indefinite \emph{by construction}, and plain CG can never run on it.
\item \textbf{The scalar $\alpha$.}  The problem's single global
  regularization parameter appears in every subdomain, so its row and column
  of $S$ are dense, wrecking both sparsity and conditioning.
\item \textbf{The cross points} --- border unknowns touched by three or more
  subdomains.  Unlike (i) these do not make $S$ indefinite; they make it
  increasingly \emph{ill-conditioned} as the subdomain count grows, which is a
  subtler defect and the subject of Section~\ref{sec:crosspoints}.
\end{enumerate}

The \emph{peel} removes both before the Krylov solve.  Split the border
index set into the peeled set $\mathrm{P}$ (all three categories above) and
the kept set $\mathrm{f}$, and partition
\begin{equation}
  \begin{bmatrix} \Sff & \SfP\\[2pt] \SfP^{\mathsf T} & \SPP \end{bmatrix}
  \begin{bmatrix} \Delta y_{\mathrm f}\\[2pt] \Delta y_{\mathrm P}\end{bmatrix}
  =
  \begin{bmatrix} r_{\mathrm f}\\[2pt] r_{\mathrm P}\end{bmatrix}.
  \label{eq:peel-system}
\end{equation}
With
\begin{equation}
  Z \;=\; \Sff^{-1}\,\SfP
  \qquad\text{and}\qquad
  T \;=\; \SPP - \SfP^{\mathsf T} Z
  \;\;\in\R^{|\mathrm P|\times|\mathrm P|}
  \label{eq:ZT}
\end{equation}
(the dense Schur complement of the split), block elimination of
\eqref{eq:peel-system} gives
\begin{align}
  \Delta y_{\mathrm P} &= T^{-1}\bigl(r_{\mathrm P}
      - \SfP^{\mathsf T}\, \Sff^{-1} r_{\mathrm f}\bigr),
  \label{eq:peel-a}\\
  \Delta y_{\mathrm f} &= \Sff^{-1} r_{\mathrm f} \;-\; Z\,\Delta y_{\mathrm P}.
  \label{eq:peel-b}
\end{align}
The indefinite and dense directions now live entirely inside $T$, and
$\Sff$ --- the operator CG actually iterates on --- is symmetric positive
definite whenever the design premise holds.  This is the FETI-DP/BDDC corner
treatment \cite{farhat2001fetidp,dohrmann2003bddc} in Schur-complement
clothing.

The size of $\mathrm P$ is set by the cross points.  On a $k\times k$ tile
partition, measured: there are $(k-1)^2$ interior cross points, each of degree
exactly three, and $4(k-1)^2$ promoted duals (four per cut-corner cell, since
the driver promotes two rank-one dual pairs there).  With $\alpha$,
\begin{equation}
  |\mathrm P| \;=\; 5(k-1)^2 + 1
  \qquad\text{(measured: } 6,\,21,\,46 \text{ at } k=2,3,4\text{)},
  \label{eq:peelsize}
\end{equation}
against interfaces of $p = 128,\,261,\,400$.  Strip partitions and 1D problems
have no cross points and no promoted duals, and there $|\mathrm P| = 1$.
Equation~\eqref{eq:peelsize} is worth reading carefully: $|\mathrm P|$ grows
\emph{linearly} with the number of subdomains $k^2$, so $T$ --- which is
precisely the coarse problem of this scheme, in the FETI-DP sense --- is only
cheap at modest $k$.  That growth is the classical scalability limit of a
coarse space solved directly, and it is where the multilevel and
inexact-coarse-solve literature applies
\cite{klawonn2007inexact,tu2007three}.  The implementation therefore carries a
guard that declines the cross-point peel outright if it would not leave
$|\mathrm P|$ a small fraction of $p$: a partition without a small coarse
space of this kind is better served by the $\alpha$ and dual peel alone than
by moving the whole interface into a dense direct solve.

\paragraph{Building the peel cache.}  Once per factorization:
\begin{enumerate}
\item $\SfP$ and $\SPP$ are read off by applying \eqref{eq:applyS} to the
  $|\mathrm P|$ unit vectors of the peeled positions --- no assembly of $S$
  is needed;
\item $Z$ is computed one \emph{CG solve per column} of $\SfP$ (every system
  involving $\Sff$ in this solver, including these, goes through CG --- the
  honest matrix-free discipline, and also the expensive part).  Each column
  is \emph{warm-started} from its last accepted value: the systems change
  slowly along the barrier path --- across a $\delta$-retry only the
  regularization moves --- and per-column telemetry showed the dominant
  failure mode was a \emph{cold-start stall} (zero progress from $x_0=0$ for
  the entire stall window, on a column that solved fine one build later), 
  which a nonzero start sidesteps by handing CG a different Krylov space;
\item $T$ is formed densely from \eqref{eq:ZT} and factorized by LU for
  \eqref{eq:peel-a};
\item $\In(T)$ is computed for the inertia prediction
  (Section~\ref{sec:inertia-pred}).
\end{enumerate}
The action of $\Sff$ on a vector is obtained for free from
\eqref{eq:applyS}: embed with zeros on the peeled positions, apply $S$, and
read back the kept positions.

If any column of $Z$ fails to converge, the cache is declared unbuildable
and the whole factorization is reported \code{SINGULAR}
(Section~\ref{sec:inertia-pred} explains why refusal is the correct
semantics).

% =============================================================================
\section{Why the cross points must be peeled}
\label{sec:crosspoints}

Categories (i) and (ii) of Section~\ref{sec:peel} are forced: without them
$\Sff$ is indefinite or dense and CG cannot run at all.  Category (iii) is
different --- $S$ restricted to the non-cross-point interface is definite, so
CG \emph{runs}; it simply converges worse and worse as the subdomain count
grows.  This section records why that happens and what it costs.

\subsection{The FETI-DP corner rule}

This is not an analogy with domain decomposition; it is the same construction.
Farhat et al.\ \cite{farhat2001fetidp} define a corner as, in their
Definition~D1, ``its cross points --- that is, the points belonging to more
than two subdomains,'' and make exactly those unknowns \emph{primal}: globally
continuous, eliminated by static condensation into a coarse problem, while the
rest of the interface is left to the iteration.  Their coarse operator
\[
  K^*_{cc} \;=\; K_{cc} \;-\; \sum_s (K^s_{rc}B^s_c)^{\mathsf T}(K^s_{rr})^{-1}(K^s_{rc}B^s_c)
\]
is precisely the $T$ of \eqref{eq:ZT}, and their reason for the choice is also
ours twice over: making the corners primal is what makes each local block
$K^s_{rr}$ nonsingular, which is exactly the rank deficiency this project's
driver cures by promoting cut-corner duals.

What makes the rule sufficient \emph{in two dimensions} is a theorem.  Mandel
and Tezaur \cite{mandel2001convergence} prove that with cross-point
constraints alone the condition number of the resulting FETI-DP operator is
bounded by $C(1+\log(H/h))^2$, \emph{independently of the number of
subdomains}.  The edge and face averages that dominate the BDDC literature are
a three-dimensional requirement --- in 3D, corners alone lose a factor $H/h$
\cite{klawonn2002dual} --- and are not needed here.  The corresponding
statement for the interface \emph{without} a coarse component is the uniform
diagnosis of the field: convergence degrades with the subdomain count because
only local information is propagated
\cite{klawonn2002dual,farhat2001fetidp}.

\subsection{The same conclusion from the optimization side}

The ASd preconditioner of Section~\ref{sec:precond} is taken from Lueg et
al.\ \cite{lueg2025}, and that paper reaches the identical conclusion about
its own method: ``it is well known that so-called one-level preconditioners as
described above do not scale well with the number of partitions $n_P$ \dots
Improved scalability is obtained by two-level methods, which add a coarse grid
correction term.''  They give the two-level form
$M_{2AS}(r) = N_0^{\mathsf T}S_0^{-1}N_0 r + \sum_k N_k S_k^{-1} N_k^{\mathsf T} r$
with $S_0 = N_0 S N_0^{\mathsf T}$, and leave its application to optimization
Schur complements as future work.

Their own PDE-constrained test problem does not exercise the difficulty: it
cuts along breakpoints in time, so that ``any set of two complicating
variables are at most shared by one partition'' --- a one-dimensional
interface with \emph{no cross points at all}, where their ASd coincides with
plain additive Schwarz.  A two-dimensional tile partition is the regime the
paper identifies as needing the second level, and peeling the cross points is
the cheapest available form of it: they are placed in $T$, which is already
built and already solved exactly, so the coarse correction costs
$(k-1)^2$ extra columns and no new machinery.

\subsection{Measured}

Cameraman, exact Hessians, cross-point peel off versus on.  The solution is
unchanged in every row (identical PSNR to two decimals):

\begin{center}
\begin{tabular}{lccc}
\toprule
run & cross points & peel off & peel on\\
\midrule
$N=64$, 4 strips        & 0 & 129 it, 7.99\,s & \emph{bit-identical}\\
$N=64$, $2\times2$ tiles & 1 & 178 it, 11.4\,s & 170 it, 10.6\,s\\
$N=32$, $3\times3$ tiles & 4 & 195 it, 6.31\,s & 113 it, 3.24\,s\\
$N=32$, $4\times4$ tiles & 9 & 2251 it, 160\,s & 509 it, 36.8\,s\\
$N=64$, $4\times4$ tiles & 9 & did not converge & 782 it, 173\,s\\
\bottomrule
\end{tabular}
\end{center}

Two features of this table are the argument.  The benefit scales with the
number of cross points, from nothing at one to a factor of four at nine; and
strip partitions, which have no cross points, are \emph{bit-identical} --- the
mechanism confirming itself rather than a general-purpose improvement that
happens to help.  The last row is the one that matters most: at $N=64$ with
$4\times4$ tiles the solver without the cross-point peel does not reach a
solution at all (it is still grinding on dual infeasibility after ten minutes,
having reached the final barrier level), and with it the run completes in
under three minutes.

% =============================================================================
\section{Preconditioned conjugate gradients on \texorpdfstring{$\Sff$}{Sff}}
\label{sec:cg}

\subsection{The ASd preconditioner}
\label{sec:precond}

A Krylov interface solve stands or falls on its preconditioner: with plain
Jacobi ($M = \diag|\Sff|$) every CG solve on a two-dimensional interface
stalled in measurement, which is why Jacobi was removed from the
implementation and only \emph{ASd} (``additive Schur, assembled diagonal'',
after equations (20)--(21) of the Lueg paper \cite{lueg2025}) remains.

For each subdomain $k$, let $R_k$ select the \emph{kept} border positions
that subdomain touches (a subset of $N_k$; the rows of $S_k$ belonging to
peeled positions are dropped).  Form the small dense block
\begin{equation}
  M_k \;=\; R_k S_k R_k^{\mathsf T}
  \quad\text{with its diagonal replaced by}\quad
  \bigl(\diag S\bigr)\big|_{R_k},
  \label{eq:asd-block}
\end{equation}
using \eqref{eq:diagS}, and set
\begin{equation}
  M^{-1} \;=\; \sum_{k=1}^{K} R_k^{\mathsf T} M_k^{-1} R_k .
  \label{eq:asd}
\end{equation}
The diagonal swap is the whole trick: $S_k$ alone describes what subdomain
$k$ does to the shared border unknowns and badly underestimates their true
self-coupling, to which every neighbouring subdomain also contributes;
$\diag(S)$ is the assembled truth, and the one quantity a distributed code
can share with a single all-reduce.  A singular local block $M_k$ is simply
skipped --- a preconditioner may be incomplete; if what remains is useless,
CG's own breakdown guard (below) catches it.

\subsection{The iteration}

\begin{algorithm}[ht]
\caption{Preconditioned CG with breakdown guards and best-iterate memory}
\label{alg:cg}
\begin{algorithmic}[1]
\State $x \gets 0$;\quad $r \gets b$;\quad $z \gets M^{-1} r$;\quad
       $\rho \gets r^{\mathsf T} z$;\quad $q \gets z$
\State $x^\star \gets 0$;\quad $\rho^\star \gets \|b\|$
       \Comment{best iterate seen, and its residual norm}
\For{$t = 1,\dots,t_{\max}$}
  \If{$\|r\| < \rho^\star$} $x^\star \gets x$, $\rho^\star \gets \|r\|$ \EndIf
  \If{$\|r\| \le \tau \|b\|$ \textbf{or} no improvement for $40$ iterations}
      \textbf{break} \EndIf
  \If{$\rho \le 0$} \textbf{break}
      \Comment{preconditioner not SPD}\EndIf
  \State $w \gets \Sff\, q$ \Comment{via \eqref{eq:applyS}, embedded}
  \If{$q^{\mathsf T} w \le 0$} flag \emph{indefinite}; \textbf{break}
      \Comment{operator not SPD along $q$}\EndIf
  \State $\alpha \gets \rho / q^{\mathsf T} w$;\quad
         $x \gets x + \alpha q$;\quad $r \gets r - \alpha w$
  \State $z \gets M^{-1} r$;\quad $\rho' \gets r^{\mathsf T} z$;\quad
         $q \gets z + (\rho'/\rho)\, q$;\quad $\rho \gets \rho'$
\EndFor
\State \Return $x^\star$ and its relative residual
\end{algorithmic}
\end{algorithm}

Algorithm~\ref{alg:cg} is textbook PCG \cite{hestenes1952} with three
provisions that turn ``CG assumes SPD'' into a runtime check.  The guard
$\rho \le 0$ certifies the preconditioner is not positive definite along the
current residual; the guard $q^{\mathsf T} \Sff q \le 0$ exhibits a direction
of non-positive curvature.  In exact arithmetic the latter is a
one-directional \emph{proof} of indefiniteness; in floating point at
$\|A\| \sim 10^{18}$ a merely tiny $q^{\mathsf T} w$ can round non-positive,
so the flag over-reports and is treated as \emph{evidence, not certificate}
(it is counted in the statistics but never acted on --- runs where it fired a
dozen times still predicted the inertia correctly on $78/79$
factorizations).  Finally the \emph{best} iterate, not the last, is returned,
so a stalled solve still hands back something usable.

\subsection{Acceptance, and what happens on rejection}
\label{sec:accept}

After the peel recombination \eqref{eq:peel-a}--\eqref{eq:peel-b} assembles
the full $\Delta y$, the answer is judged by the residual of the \emph{full}
interface system --- not the $\Sff$ subproblem CG actually saw, because the
peel algebra sits between the two:
\[
  \frac{\|r_y - S\,\Delta y\|}{\|r_y\|} \;<\; 10^{-2}
  \qquad\text{and}\qquad
  \text{(CG's own relative residual)} \;<\; 10^{-2}.
\]
A rejected answer has nowhere to fall back to --- $S$ was never factorized
--- so it is handed back \emph{anyway}, as the best iterate available, and
the outer iterative refinement (Section~\ref{sec:refine}) judges it against
the true matrix.  Reporting \code{SINGULAR} instead would be truthful but
useless: IPOPT would answer with $\delta_w$, which cannot repair a Krylov
convergence failure, and the run dies (measured, in the historical
no-inertia mode: internal error on three of four configurations).  Only a
solve that produced no iterate at all is reported as failed.

\subsection{Relation to the iterative strategy of Lueg et al., \S 3.1}
\label{sec:lueg-audit}

The interface solve of this solver is an implementation of the iterative
strategy of \cite{lueg2025}, \S 3.1 --- PCG on the Schur complement system,
never forming or factorizing $S$, with a mandatory preconditioner, local step
recovery by parallel subdomain back-solves (their eq.~11 $=$ our
\eqref{eq:solve-c}), local inertia read from the pivot signs of the $W_k$
factorizations, and negative-curvature monitoring inside PCG --- with two
departures, one deliberate and one forced, and two points on which this
implementation realizes what their paper lists as future extensions.

\paragraph{Deliberate departure: how $S$ is applied.}
Their eq.~(12) applies $S\cdot u$ matrix-free, at one $W_k$ back-solve per
partition per CG iteration, so that not even the local Schur blocks exist.
Here the $S_k$ are formed once per factorization (Section~\ref{sec:sk}) and
each application is a dense $p_k\times p_k$ product.  The operator and the
one-reduction communication pattern are identical; only the local cost model
differs, and the measurements of Section~\ref{sec:sk} decide it: the
matrix-free application costs roughly ten times the dense product per
iteration, and at the observed $10^2$--$10^3$ CG iterations per factorization
the formation cost is repaid twenty-fold.  Two further facts tilt the trade.
Their own preconditioners (eqs.~17 and 20--21) require the local Schur blocks
regardless, and their cost model for formation --- $p_k$ full back-solves per
partition --- is exactly what the forward-only, static-reach route of
Section~\ref{sec:sk} removes.

\paragraph{Forced departure: the peel.}
Their premise ``$S$ is positive definite through inertia correction'' is a
structural fact of their formulation: the complicating variables are
primal-only and appear only in linear linking constraints, so no dual ever
reaches their border (and their corner block is zero, where our
pure-permutation $C$ is not).  In the present formulation the border carries
promoted corner duals --- $S$ is indefinite by construction --- plus $\alpha$
and the cross points, so \S 3.1 is inapplicable to $S$ itself.  The peel of
Section~\ref{sec:peel} is the adaptation that restores their premise:
$\Sff$ plays the role their $S$ plays, and CG runs there.

\paragraph{Inertia checks.}
For the iterative mode their \S 3.5 offers two ways to check
$\In(S) = (p,0,0)$: monitoring $p^{\mathsf T} S p \ge 0$ inside PCG (``this
is not a rigorous check,'' their words), and a conservative per-block test
$\In(S_k) = (p_k, 0, 0)$.  The first is precisely the $pAp$ guard of
Algorithm~\ref{alg:cg}, retained here as guard and telemetry.  The second is
inapplicable in this setting: our $S_k$ carry dual contributions and are
indefinite by design.  The prediction of Section~\ref{sec:inertia-pred}
replaces both with an exact count.

\paragraph{Their future work, implemented here.}
Their implementation notes record that the absence of a filter line search
and of iterative refinement forced PCG to a $10^{-9}$ relative residual, name
both as ``sensible extensions,'' and flag the interplay of iterative
tolerance with step quality as future research.  This solver has both --- it
sits inside IPOPT's filter line search, and Section~\ref{sec:refine} is the
refinement --- which is what lets the acceptance threshold sit at $10^{-2}$;
the tolerance sweep reported in Section~\ref{sec:measured} is the experiment
their future-work paragraph describes.  One element of their \S 3.5 is
genuinely not adopted: \emph{localized} regularization
($\delta^k_H, \delta^k_C$ adjusted only in the offending partition) is
unavailable to a solver behind IPOPT's linear-solver interface, where the
regularization arrives already applied globally and the contract is limited
to reporting \code{SINGULAR} or wrong inertia.

\subsection{Removing the forced departure: the consensus formulation}
\label{sec:consensus}

The peel is an adaptation to a premise failure, and the failure is
formulation-induced --- so it can be removed at the source.  The companion
class \code{mpcc\_2d\_consensus\_tnlp.hpp} rebuilds the same MPCC in Lueg's
duplicate-and-link form: every constraint row is assigned wholly to one tile,
every variable referenced by rows of two or more tiles (the border $u$ nodes,
the cut-line $q$ components, and $\alpha$ --- never $r,\delta,\theta$, which
are strictly cell-local) receives one local copy per referencing tile, and
the original index becomes a consensus variable appearing only in the linear
linking rows $v_k - v = 0$.  The reformulation is exact: IPOPT's monolithic
solver run on both formulations agrees on $\alpha^\star$ to six digits.  The
solver of this paper is untouched --- it consumes the new problem through the
same injected owner map, and its partition-leak check doubles as an audit of
the construction.

Structurally, the premise of \S 3.1 is thereby restored: rows are never cut,
so the rank deficiencies that forced corner-dual promotion cannot arise, no
dual reaches the border, and $S$ is positive definite once inertia correction
succeeds; $\alpha$'s dense row collapses to $K$ linking couplings.  The dual
peel becomes vacuous, while the cross points survive as consensus variables
of degree three and are still handled by the (now purely primal) peel.
Measured, with the cross-point peel active in both columns:

\begin{center}
\begin{tabular}{lcc}
\toprule
run & permutation & consensus\\
\midrule
$N=16$, $2\times2$ tiles & 0.19\,s, 53 it & 0.13\,s, 56 it\\
$N=32$, $3\times3$ tiles & 3.16\,s, 113 it & 1.26\,s, 80 it\\
$N=32$, $4\times4$ tiles & 35.5\,s, 509 it & \textbf{2.06\,s, 84 it}\\
$N=64$, $2\times2$ tiles & 10.6\,s, 170 it & 6.60\,s, 139 it\\
$N=64$, $4\times4$ tiles & 172\,s, 782 it & \textbf{13.0\,s, 128 it}\\
$N=64$, $4$ strips & 8.05\,s, 129 it & 8.05\,s, 136 it\\
\bottomrule
\end{tabular}
\end{center}

Identical reconstructions in every row.  The headline is the flatness of the
consensus column in the tile count ($80 \to 84$ at $N=32$, $139 \to 128$ at
$N=64$, against $113 \to 509$ and $170 \to 782$ for the permutation form):
with the formulation-induced indefiniteness gone, the decomposition scales
the way the theory of Section~\ref{sec:crosspoints} says it should.  The
price is a larger NLP (one copy per shared variable per tile plus as many
linking rows, $O(kN)$), a different barrier trajectory --- runs match the
monolithic reference in solution, not iteration-for-iteration --- and one
benign departure from strict Lueg form: the objective stays on the consensus
variables, contributing a positive semidefinite diagonal to $C$, which can
only help the definiteness of $S$.

% =============================================================================
\section{The predicted inertia}
\label{sec:inertia-pred}

\subsection{Haynsworth, a second time}

The classical objection to a Krylov interface solve inside IPOPT is that the
inertia \eqref{eq:haynsworth} still needs $\In(S)$, which classically comes
only from a factorization of $S$ --- so $S$ would be assembled and factorized
anyway, the direct back-solve would then be nearly free, and CG would be
strictly overhead.  The resolution is to \emph{derive} $\In(S)$ from
something already at hand.  Theorem~\ref{thm:haynsworth} applies to the peel
split \eqref{eq:peel-system} exactly as it applied to the arrowhead:
\begin{equation}
  \In(S) \;=\; \In(\Sff) \;+\; \In(T),
  \qquad T = \SPP - \SfP^{\mathsf T}\, \Sff^{-1}\, \SfP,
  \label{eq:haynsworth2}
\end{equation}
and $T$ is \emph{already built} --- it is the peel cache of
Section~\ref{sec:peel}, of size $|\mathrm P| \le$ a few dozen, so $\In(T)$
costs a dense symmetric eigendecomposition of a tiny matrix: nothing.
The entire prediction is therefore the single assumption
\begin{equation}
  \In(\Sff) \;=\; (\,|\mathrm f|,\,0,\,0\,),
  \qquad\text{i.e.\ } \Sff \succ 0.
  \label{eq:premise}
\end{equation}
What makes \eqref{eq:premise} the right assumption to bet on is that it is
\emph{also exactly what CG requires}: the peel exists precisely to make
$\Sff$ definite.  The scheme is self-consistent --- the interface operator is
usable exactly when the reported inertia is right --- and the reported total
is
\begin{equation}
  n_-(A) \;=\; \sum_{k=1}^K n_-(W_k) \;+\; n_-(T).
  \label{eq:reported}
\end{equation}

Three implementation details keep \eqref{eq:reported} trustworthy.  First,
$\In(T)$ is computed by an exact dense \emph{symmetric eigendecomposition},
deliberately not by a dense $LDL^{\mathsf T}$: Eigen's dense \code{LDLT} is a
pivoted Cholesky for semidefinite matrices, not Bunch--Kaufman, and its pivot
signs are unreliable on indefinite input (the production solver measured
$485$--$491$ reported negatives where the truth was $512$).  Second, $T$ is
diagonally \emph{equilibrated} before the eigendecomposition:
$T \to DTD$ with $D = \diag(|T_{ii}|^{-1/2})$ is a congruence, so
$\In(DTD) = \In(T)$ by Sylvester's law, while the eigenvalue scales become
comparable --- necessary because $T$ mixes the $\alpha$ direction
(barrier-scaled, up to $\sim\!10^{20}$) with corner-dual directions
($\sim\!10^{-9}$), and a zero test against the \emph{global} eigenvalue scale
mistakes honest tiny eigenvalues for noise (measured: at $N=64$, $4\times4$
tiles it refused every factorization this way, and the run could not start;
a referee $LDL^{\mathsf T}$ of the assembled $S$ scored the equilibrated
test's predictions $199/200$ correct).  Third, an eigenvalue of the
equilibrated matrix with $|\lambda| < 10^{-12}\max_i|\lambda_i|$ still means
$\In(T)$ is not well determined, and the prediction is refused.

\subsection{Refusing to predict}

If the peel cache cannot be built --- a CG solve for a column of $Z$ fails
to converge, or $T$ has an undeterminable eigenvalue --- \emph{no prediction
is issued and the factorization is reported \code{SINGULAR}}.  This is the
safety property of the whole design: a prediction the solver cannot stand
behind would corrupt IPOPT's $\delta_w$ loop silently, whereas
\code{SINGULAR} merely costs a re-factorization at a larger $\delta_w$ ---
a failure IPOPT knows how to cure.

\subsection{Why \texorpdfstring{$n_-(S)$}{n-(S)} is predictable at all}
\label{sec:why-predictable}

Write $n_-(A) = m + \nu$, where $m$ is the number of dual unknowns and $\nu$
the negative-curvature excess of the reduced Hessian --- the quantity IPOPT
drives to zero with $\delta_w$.  Each $W_k$ is itself a saddle-point matrix,
so $n_-(W_k) = m_k + \nu_k$ with $m_k$ its interior duals.  Every dual row is
either owned by one subdomain or promoted to the border, so
$\sum_k m_k = m - p_{\mathrm{dual}}$, where $p_{\mathrm{dual}}$ is the
number of promoted (peeled) border duals.  Substituting into
\eqref{eq:haynsworth},
\begin{equation}
  n_-(S) \;=\; p_{\mathrm{dual}} \;+\; \nu \;-\; \sum_{k=1}^{K}\nu_k ,
  \label{eq:nu-count}
\end{equation}
which explains the measured fact that $n_-(S)$ equals the promoted-dual count
whenever the local and global curvature counts agree
($\nu = \sum_k \nu_k$, e.g.\ both zero at accepted steps).  The
$\In(T)$ route \eqref{eq:haynsworth2} is used in practice because it needs no
such assumption about $\nu$: $T$ sees whatever the corner duals actually
contribute.

% =============================================================================
\section{Iterative refinement and the outer solve}
\label{sec:refine}

Every solve of the full system runs the arrowhead recursion
\eqref{eq:solve-a}--\eqref{eq:solve-c} inside a few sweeps of iterative
refinement measured against the \emph{original input triplets}:
\begin{equation}
  \text{repeat up to 3 times:}\quad
  \rho = b - A x
  \;\;\text{(triplet matvec)},\qquad
  x \gets x + \mathcal{A}^{-1}\rho,
  \label{eq:refine}
\end{equation}
where $\mathcal{A}^{-1}$ denotes one arrowhead solve, stopping early at
relative residual $10^{-11}$ and always returning the \emph{best} iterate
seen.  This is not a luxury.  Near-singular pivots --- IPOPT's own
$\delta_c = 10^{-8}\mu^{1/4}$ on the dual directions, for instance --- lose
about ten digits through the Schur recursion (measured: relative residual
$\sim\!0.1$ unrefined against $\sim\!10^{-12}$ refined).  Because the
residual in \eqref{eq:refine} is evaluated straight from the input triplets,
refinement checks the \emph{decomposition itself}, not merely the arithmetic
inside it; it is also the safety net that absorbs the accepted-anyway CG
iterates of Section~\ref{sec:accept}.  Returning one's best answer at a nasty
iterate and letting the globalization cope is exactly what a monolithic
sparse solver does too.  The solve reports failure upward only if the
residual never once evaluated finite.

\subsection{Life cycle}

\begin{center}
\begin{tabular}{lp{10.2cm}}
\toprule
\code{set\_structure} & once per sparsity pattern: number the unknowns,
  route every triplet to its block ($W_k$, $B_k$, or $C$), detect partition
  leaks \eqref{eq:noleak}, build the peel sets, run the symbolic analysis of
  every $W_k$.\\[3pt]
\code{factorize} & every Newton step: (1) factorize all $W_k$
  \eqref{eq:ldlt}, accumulate $\sum_k \In(W_k)$ --- any breakdown
  $\Rightarrow$ \code{SINGULAR}; (2) form all $S_k$ by
  \eqref{eq:forward-only}--\eqref{eq:row-contraction}; (3) assemble $C$ and
  $\diag(S)$ \eqref{eq:diagS}; (4) build the peel cache
  ($\SfP$, $\SPP$, $Z$, $T$, $\In(T)$) and the ASd blocks --- failure
  $\Rightarrow$ refusal $\Rightarrow$ \code{SINGULAR}; report
  \eqref{eq:reported}.  Steps 1--2 are independent per subdomain.\\[3pt]
\code{solve} & per right-hand side: refinement loop \eqref{eq:refine} around
  the arrowhead recursion; the interface step \eqref{eq:solve-b} is
  Algorithm~\ref{alg:cg} plus the peel recombination
  \eqref{eq:peel-a}--\eqref{eq:peel-b}.\\
\bottomrule
\end{tabular}
\end{center}

\subsection{Constants}

\begin{center}
\begin{tabular}{lll}
\toprule
quantity & value & role\\
\midrule
\code{cg\_tol} & $10^{-10}$ & CG target relative residual\\
\code{cg\_maxit} & $500$ & CG iteration cap\\
stall window & $40$ & iterations without improvement before giving up\\
acceptance & $10^{-2}$ & max relative residual (CG and full interface)\\
peel-column acceptance & $10^{-2}$ & max relative residual for a $Z$ column\\
refinement target & $10^{-11}$ & early exit of \eqref{eq:refine}\\
refinement sweeps & $3$ & maximum corrections\\
$\In(T)$ threshold & $10^{-12}\,\|\lambda\|_{\max}$ & below: refuse to
  predict\\
\bottomrule
\end{tabular}
\end{center}

% =============================================================================
\section{What the measurements decided}
\label{sec:measured}

The single configuration described above --- CG-only interface, ASd-only
preconditioning, predicted-only inertia --- is the surviving point of a
larger design space measured in this codebase (the production solver
\code{dd\_solver.hpp} retains the rest: MA57/MA97/MUMPS backends, a direct
and a MINRES interface route, lagged Schur caching, a nested second level).
The decisive measurements, all on the cameraman image with exact Hessians:

\paragraph{Correctness against a monolithic reference.}  Against IPOPT's own
MUMPS on identical instances, the decomposition changes the answer not at
all: identical $\alpha^\star$ to six digits and identical PSNR at every
size/partition tested; the $\sim\!15\%$ extra IPOPT iterations are the
unpivoted $LDL^{\mathsf T}$ asking for regularization MA57 would not have
needed (Section~\ref{sec:ldlt}).

\paragraph{The preconditioner is everything.}  With ASd, CG carried every
interface solve on strips and nearly all on tiles; with Jacobi on a
$2\times 2$ tile interface, CG carried \emph{none} ($0/50$).  Hence Jacobi's
removal.

\paragraph{The three inertia modes} (historical; \emph{exact} assembled and
factorized $S$, \emph{none} reported no inertia and used the curvature test
of \cite{chiang2016inertia}, \emph{predicted} is Section~\ref{sec:inertia-pred}):
\begin{center}
\begin{tabular}{lccc}
\toprule
run & exact & predicted & none\\
\midrule
$N=16$, 2 strips        & 50 it, 0.11\,s & 50 it, 0.09\,s & 71 it, 0.11\,s\\
$N=16$, $2\times2$ tiles & 50 it, 0.10\,s & 80 it, 0.21\,s & 108 it, 0.24\,s\\
$N=32$, 3 strips        & 75 it, 1.49\,s & 80 it, 1.48\,s & 115 it, 1.83\,s\\
$N=32$, $3\times3$ tiles & 75 it, 1.72\,s & 125 it, 4.51\,s & 1051 it, 63.3\,s\\
\bottomrule
\end{tabular}
\end{center}
All modes reach the same solution (PSNR identical to two decimals).  Scored
against a referee factorization of $S$, the prediction was correct on
$78/79$, $99/99$, $111/117$ and $132/133$ factorizations respectively, and
every mismatch was off by exactly one.  On strips the prediction is
\emph{free} --- same iterations and wall clock as exact inertia, with the one
serial factorization gone --- which is the result the design was after.  What
still costs on tiles is not wrong predictions but \emph{refusals}: peel-cache
columns whose CG solve does not converge.  At $N=64$ with $4\times4$ tiles the
refusals become total (every factorization refused; the run cannot start), so
the remaining obstacle of the approach is preconditioning the peel-column
solves --- a tractable, distributed-friendly problem --- and no longer ``we
must factorize $S$ to know its inertia.''

\paragraph{Formation of $S_k$.}  Forward-only + sparse-RHS
(Sections \ref{sec:sk}--\ref{sec:sparse-fwd}) gave $1.65$--$1.8\times$
whole-run speedups over the naive route at $N=32$--$64$; the row contraction
(Section~\ref{sec:row-contraction}) a further $1.1$--$2.6\times$.

\paragraph{Validation gates.}  A standalone smoke test
(\code{dd\_simple\_smoke.cpp}, pure Eigen, no IPOPT) checks the arrowhead
solve against a dense LU and the predicted inertia \eqref{eq:reported}
against a dense symmetric eigendecomposition, on synthetic arrowhead KKT
matrices with and without the peel.  \code{DDS\_SCHUR\_CHECK=1} referees
every $S_k$ at runtime; \code{DDS\_DEBUG=1} prints partition and CG
diagnostics.

% =============================================================================
\begin{thebibliography}{9}

\bibitem{waechter2006implementation}
A.~W\"achter and L.~T. Biegler,
\emph{On the implementation of an interior-point filter line-search
algorithm for large-scale nonlinear programming},
Mathematical Programming 106 (2006) 25--57.

\bibitem{haynsworth1968}
E.~V. Haynsworth,
\emph{Determination of the inertia of a partitioned Hermitian matrix},
Linear Algebra and its Applications 1 (1968) 73--81.

\bibitem{vanderbei1995symmetric}
R.~J. Vanderbei,
\emph{Symmetric quasidefinite matrices},
SIAM Journal on Optimization 5 (1995) 100--113.

\bibitem{chiang2016inertia}
N.-Y. Chiang and V.~M. Zavala,
\emph{An inertia-free filter line-search algorithm for large-scale nonlinear
programming},
Computational Optimization and Applications 64 (2016) 327--354.

\bibitem{gilbert1988}
J.~R. Gilbert and T.~Peierls,
\emph{Sparse partial pivoting in time proportional to arithmetic
operations},
SIAM Journal on Scientific and Statistical Computing 9 (1988) 862--874.

\bibitem{mandel2001convergence}
J.~Mandel and R.~Tezaur,
\emph{On the convergence of a dual-primal substructuring method},
Numerische Mathematik 88 (2001) 543--558.

\bibitem{klawonn2002dual}
A.~Klawonn, O.~B. Widlund, and M.~Dryja,
\emph{Dual-primal FETI methods for three-dimensional elliptic problems with
heterogeneous coefficients},
SIAM Journal on Numerical Analysis 40 (2002) 159--179.

\bibitem{klawonn2007inexact}
A.~Klawonn and O.~Rheinbach,
\emph{Inexact FETI-DP methods},
International Journal for Numerical Methods in Engineering 69 (2007)
284--307.

\bibitem{tu2007three}
X.~Tu,
\emph{Three-level BDDC in three dimensions},
SIAM Journal on Scientific Computing 29 (2007) 1759--1780.

\bibitem{hestenes1952}
M.~R. Hestenes and E.~Stiefel,
\emph{Methods of conjugate gradients for solving linear systems},
Journal of Research of the National Bureau of Standards 49 (1952) 409--436.

\bibitem{farhat2001fetidp}
C.~Farhat, M.~Lesoinne, P.~LeTallec, K.~Pierson, and D.~Rixen,
\emph{FETI-DP: a dual--primal unified FETI method I},
International Journal for Numerical Methods in Engineering 50 (2001)
1523--1544.

\bibitem{dohrmann2003bddc}
C.~R. Dohrmann,
\emph{A preconditioner for substructuring based on constrained energy
minimization},
SIAM Journal on Scientific Computing 25 (2003) 246--258.

\bibitem{lueg2025}
L.~R. Lueg, M.~L. Bynum, C.~D. Laird, and L.~T. Biegler,
\emph{Domain decomposition preconditioners for Schur complement systems
arising in structured nonlinear optimization problems},
Optimization and Engineering 27 (2026) 555--585
(published online 2 September 2025), \texttt{doi:10.1007/s11081-025-10020-1}.
The ASd preconditioner used here is their equation (21); the two-level
correction discussed in Section~\ref{sec:crosspoints} is their equation (22).

\end{thebibliography}

\end{document}
